\documentclass[11pt]{article}

% Course-project Phase-1 proposal (Research track). 1-2 pages; refs excluded.
\usepackage{acl}          % final mode (no line numbers)
\usepackage{times}
\usepackage{latexsym}
\usepackage[T1]{fontenc}
\usepackage[utf8]{inputenc}
\usepackage{microtype}
\usepackage{inconsolata}
\usepackage{amsmath,amssymb}

\title{The Reliability Tax: When Thinking Makes Language Models Less Trustworthy}
\author{Anonymous \\
  Natural Language Processing, Course Project (Research Track) \\
  \texttt{anonymous@example.com}}

\begin{document}
\maketitle

\begin{abstract}
Reasoning-augmented language models are deployed by default for high-stakes
queries, and how long a model ``thinks'' is now a user-facing knob. The
prevailing assumption is that more thinking yields more trustworthy answers.
We test it. We formalize the capability--reliability trade-off as a measurable
\emph{reliability tax}, and ask whether it is a law---and whether it is
\emph{predictable}. A black-box budget sweep over four task types measures
accuracy and three reliability axes jointly (calibration, selective
prediction, abstention-versus-fabrication). A first real run already shows the
tax is \emph{task-dependent} and \emph{switch-triggered}: on false-premise
items, enabling thinking raises fabrication from 3.8\% to 21.3\%, while
mathematical reasoning is unaffected.
\end{abstract}

\section{Background \& Motivation}
Reasoning-augmented language models (LRMs) spend a variable amount of compute
``thinking'' before answering \citep{Guo2025DeepSeekR1}, and this budget is
increasingly exposed as a user-facing control \citep{Muennighoff2025s1,
Snell2024TestTimeCompute}. The implicit contract is simple: \emph{more thinking
buys more trustworthy answers}. Practitioners act on it---raising the budget to
``max'' whenever an answer matters.

That contract is being contradicted from several directions.
Reinforcement learning that raises accuracy degrades calibration
\citep{Wang2026CAPO, Jin2025Entropy}; enhancing reasoning amplifies tool
hallucination, even at inference time \citep{Yin2025ReasoningTrap}; ``slower''
models are measurably \emph{less} truthful under misleading input
\citep{Fang2025Slower}; and reasoning models stay confidently wrong on a large
fraction of false-premise queries \citep{Sathyanathan2025Presuppositions}.
Crucially, the evidence is \emph{single-axis}, \emph{training-side}, and
\emph{descriptive}: it shows the trade-off \emph{exists}, not whether it is a
law, or---the property a deployment rule needs---whether it is
\emph{predictable} from the trace before a budget is committed.

\section{Research Questions \& Hypotheses}
\textbf{RQ.} Does increasing a model's thinking budget systematically trade
reliability for accuracy, and can the degradation be foreseen?
\begin{itemize}\itemsep1pt
  \item \textbf{H1 (task-dependence).} The ratio of reliability loss to
  accuracy gain differs across task types by at least a factor of two.
  \item \textbf{H2 (predictability, primary).} A predictor trained on a
  held-out half decides, at low budget and \emph{without} seeing the
  high-budget answer, whether continuing will help; the induced policy beats
  the best fixed budget in \emph{usable accuracy} at matched cost.
  \item \textbf{H3 (negative control).} Shuffling the trace features removes
  H2's advantage---if not, the predictor exploits an artefact.
\end{itemize}
Each claim carries an explicit falsification condition, and the design is
pre-registered on a held-out prompt split.

\section{Preliminary Literature}
The closest work varies the reasoning budget and measures \emph{one} axis:
\citet{DontThinkTwice2025} show over-reasoning impairs confidence calibration;
\citet{WhenMoreThinkingHurts2026} characterize the diminishing---sometimes
negative---marginal utility of extra tokens. \citet{Yin2025ReasoningTrap} prove
the capability--reliability trade-off is method-agnostic, and
\citet{Fang2025Slower} establish inverse truthfulness scaling.
Our study differs on three points: it measures three reliability axes
\emph{jointly} to test whether they degrade as one; it spans tasks and model
families rather than one benchmark; and it asks whether the degradation is
\emph{predictable}, turning the trade-off curve into a black-box budget rule.
(40 references collected and verified; ACL~2026 landscape survey of 1{,}409
papers informs the framing.)

\section{Method \& Experimental Design}
\paragraph{The object.} Everything is read off the \emph{budget--reliability
profile} of a model on a task: accuracy and each reliability axis as a function
of the thinking budget. We report against the \emph{realized} reasoning-token
count, since nominal budget knobs are not comparable across providers.

\paragraph{The tax.} With every reliability score oriented so that higher is
better, $\tau_r=(\partial\rho_r/\partial b)/(\partial \mathrm{acc}/\partial b)$
is the reliability given up per unit of accuracy gained. Because $\tau_r$
diverges where accuracy saturates, our primary summaries are bounded instead:
the fixed-coverage risk $R(c)$ and the area under the risk--coverage curve
\citep{Geifman2017Selective}.

\paragraph{Three axes.} (i) \emph{Calibration}: ECE \citep{Guo2017ECE} and
verbalized overconfidence \citep{Kadavath2022KnowWhatTheyKnow,Tian2023JustAsk}.
(ii) \emph{Selective prediction}: $R(c)$ and AURC. (iii)
\emph{Abstention vs.\ fabrication} on false-premise items, classified as
refute / hedge / endorse, with fabrication rate = endorse share
\citep{Sathyanathan2025Presuppositions,Kirichenko2025AbstentionBench}.

\paragraph{Tasks and budget.} Four task types span where compute should help
and where it can hurt: mathematical reasoning (GSM8K, MATH-500; verifiable),
factual QA (NQ), multi-hop QA (HotpotQA; composition), and false-premise items
built from SciFact claims whose evidence contradicts them
\citep{Wadden2020SciFact}. The budget axis uses the provider's official
controls---thinking on/off and reasoning-effort levels---not \texttt{max\_tokens},
which would confound budget with truncation; realized reasoning tokens are
logged per response. Four budgets $\times$ two prompt templates $\times$
$k{=}4$ samples per item, with prompt-clustered bootstrap intervals.

\paragraph{From profile to decision.} The profile yields an oracle: the budget
maximizing \emph{usable accuracy}, i.e.\ accuracy at the coverage where ECE
stays below a target $\epsilon$. We then test whether that oracle is
\emph{findable} from low-budget features alone (realized length, verbalized
confidence, self-consistency, hedging markers), with the H3 permutation
control.

\section{Preliminary Result}
A first real run on \texttt{deepseek-flash} (1{,}280 calls, 0 errors;
${\sim}\text{\textyen}7$--13) already indicates the tax is
\emph{task-dependent} and \emph{switch-triggered}:
\begin{center}\small
\begin{tabular}{lcccc}
\hline
Task & Acc.\ off$\to$max & ECE off$\to$max & Fabrication off$\to$max \\
\hline
Math          & .988$\to$.950 & .01$\to$.04 & --- \\
Factual QA    & .475$\to$.550 & .45$\to$.39 & --- \\
Multi-hop QA  & .463$\to$.562 & .46$\to$.41 & --- \\
False-premise & .962$\to$.800 & .03$\to$.13 & \textbf{.038$\to$.213} \\
\hline
\end{tabular}
\end{center}
On false-premise items---where there is \emph{no} accuracy to gain---enabling
thinking multiplies fabrication by 5.6$\times$ and breaks calibration, while
mathematical reasoning is untouched. Notably the effect is a \emph{step} at the
thinking on/off boundary rather than a smooth function of effort level, which
sharpens H1 into a stronger, more falsifiable claim.

\section{Expected Outcomes}
(1) A cross-task, cross-family characterization of whether the trade-off is a
law or a set of local effects, and where it is worst. (2) A held-out,
cost-matched evaluation of an oracle-free \emph{stopping rule} that decides
when to stop thinking. (3) A released protocol, prompts, and raw responses for
reproducibility.

\section{Timeline \& Responsibilities}
Single-member project (all roles by the author).
\begin{itemize}\itemsep1pt
  \item \textbf{Wk 1--2} (done): landscape survey; protocol and code built;
  pipeline validated; first real run on DeepSeek.
  \item \textbf{Wk 3}: scale DeepSeek to 200 items $\times\,k{=}4$; add Qwen3,
  GLM, Doubao for cross-family coverage.
  \item \textbf{Wk 4}: train and evaluate the stopping rule (H2/H3); statistics.
  \item \textbf{Wk 5}: final report, slides, and artifact release.
\end{itemize}

\bibliography{refs}

\end{document}
